\documentclass[11pt]{article}
\usepackage{amsmath,amssymb}
\usepackage[margin=1in]{geometry}

\newcommand{\Qbar}{\overline Q}

\begin{document}

\section*{Special one-coordinate limits of the two-site wavefunction}

Consider the spin-$S/2$ two-site XXX chain in the companion frame,
with
\[
\chi=\lambda+\lambda^{-1},
\qquad
Q(u)=\lambda^{u/h}q_M(u),
\qquad
\deg q_M=M.
\]
Assume that $\lambda\neq\pm1$ and set
\[
R=\lambda,
\qquad
r=\lambda^{-1}.
\]
These are the two fixed points of the projective companion twist.

\subsection*{Twist-weight basis}

Define
\[
\phi_m(z)=(z-R)^m(z-r)^{S-m},
\qquad 0\leq m\leq S.
\]
The companion-twist charge is diagonal in this basis.  Consequently,
the wavefunction belonging to a Baxter polynomial of degree $M$ obeys
the selection rule
\[
\Psi_M(z_1,z_2)
=
\sum_{\substack{0\leq a,b\leq S\\a+b=M}}
C_{ab}\,\phi_a(z_1)\phi_b(z_2).
\]

\subsection*{The specialization $z_2=\lambda$}

At $z_2=R$, only $\phi_0(R)$ is nonzero.  Therefore
\[
\boxed{
\Psi_M(z,R)
=
\begin{cases}
C_{M0}(R-r)^S\phi_M(z),&M\leq S,\\[2mm]
0,&M>S.
\end{cases}}
\]
Thus, for $M\leq S$, the complete two-site wavefunction collapses to
the single known polynomial
\[
\phi_M(z)=(z-\lambda)^M(z-\lambda^{-1})^{S-M}.
\]

In the standard monic algebraic-Bethe-ansatz normalization,
\[
C_{M0}
=
(-h)^M\frac{S!}{(S-M)!}\,q_M(\theta_2).
\]
Hence
\[
\boxed{
\Psi_M(z,\lambda)
=
(\lambda-\lambda^{-1})^S
(-h)^M\frac{S!}{(S-M)!}\,
q_M(\theta_2)
(z-\lambda)^M(z-\lambda^{-1})^{S-M}}
\]
for $M\leq S$.  A different normalization of the full eigenvector
only changes the overall scalar.

This scalar is itself a single contour bracket.  With
\[
\Qbar(u)=Q(u-Sh),
\qquad
\Pi_S(u)=\prod_{k=0}^{S-1}Q_\theta(u-kh),
\]
one has
\[
q_M(\theta_2)
=
\lambda^{-\theta_2/h}
\frac{\langle\Pi_S(u)\Qbar(u)\rangle_2}
     {\langle\Pi_S(u)\rangle_2}.
\]

\subsection*{The specialization $z_2=\lambda^{-1}$}

At $z_2=r$, only $\phi_S(r)$ is nonzero.  Thus
\[
\boxed{
\Psi_M(z,r)
=
\begin{cases}
0,&M<S,\\[2mm]
C_{M-S,S}(r-R)^S\phi_{M-S}(z),&M\geq S.
\end{cases}}
\]
Equivalently, for $M\geq S$ its $z$-dependence is
\[
\Psi_M(z,\lambda^{-1})
\propto
(z-\lambda)^{M-S}
(z-\lambda^{-1})^{2S-M}.
\]
The coefficient is most naturally expressed through the dual, or hole,
Baxter polynomial of degree $2S-M$.

This point has the additional simplification that the Baxter exponential
cancels inside the second bracket:
\[
r^{(u-\theta_2)/h}\Qbar(u)
=
\lambda^{\theta_2/h-S}q_M(u-Sh).
\]

\subsection*{A sector-independent projective limit}

If one wants a generically nonzero simplification in every sector, the
best point is projective infinity:
\[
\Psi_M(z,\infty)
:=
\lim_{w\to\infty}w^{-S}\Psi_M(z,w).
\]
Let
\[
d=\frac{\theta_1-\theta_2}{h},
\qquad
\delta_S(v)
=
\frac{1}{h^S S!}
\prod_{k=0}^{S-1}(v-\theta_2-kh),
\]
and
\[
P_2(u)=\prod_{k=0}^{S}(u-\theta_2-kh).
\]
Then the higher-spin factorizing twist gives
\[
\boxed{
\Psi_M(z,\infty)
=
-\frac{h^{-S}}{(d+1)_S}
\left\langle\delta_S(v)\Qbar(v)\right\rangle_2
\left\langle
P_2(u)z^{(u-\theta_1)/h}\Qbar(u)
\right\rangle_1.}
\]
The first factor selects the endpoint $v=\theta_2+Sh$, where
$\Qbar(v)=Q(\theta_2)$.  The insertion $P_2(u)$ cancels the complete
site-$2$ string in the first measure, leaving a genuine length-one
Mellin--Baxter transform.

The useful choices are therefore
\begin{itemize}
\item $z_0=\lambda$ for sectors $M\leq S$;
\item $z_0=\lambda^{-1}$ for sectors $M\geq S$;
\item $z_0=\infty$ for a sector-independent product of two
      one-dimensional brackets.
\end{itemize}
The finite point $z_0=0$ is only triangular in the present
$\mathcal F_{12}$ ordering and does not factorize completely.

\end{document}
